\chapter{REVISIÓN CRÍTICA DE LA LITERATURA}

Este capítulo revisa las fuentes que definen el estado del arte relevante para CAVA. Empieza
por la línea de resumen de video multimodal de la que proviene este trabajo (A2Summ y
A2Summ-X), donde se originó el diagnóstico contra la fusión. Sigue con la segmentación de
escenas y capítulos (BaSSL, Scene-VLM, VidChapters), la segmentación temática de texto, la
discusión entre fusión temprana y tardía y los ensambles de agrupamientos, que son el
antecedente formal más cercano del consenso de votos. Cierra con el sistema previo SegLoop y
con el uso de VLM como jueces. Cada sección analiza los resultados reportados por sus autores
y la última resume el aporte de la tesis.

\section{Resumen de video multimodal: A2Summ}

He et al.\ \cite{he2023a2summ} proponen A2Summ, un transformer que alinea las secuencias de
video y texto y las fusiona con dos pérdidas contrastivas, evaluado en SumMe
\cite{gygli2014summe}, TVSum \cite{song2015tvsum} y conjuntos de resumen multimodal con
salida textual.

\subsection{Análisis de resultados}
El artículo reporta F1 $\approx 55.0$ en SumMe y $\approx 63.0$ en TVSum, por encima de
métodos previos basados en dispersión o grafos \cite{chu2018spell}. El grupo replicó el
pipeline completo y obtuvo los mismos números que el checkpoint oficial (F1 0.625 en TVSum;
0.6383 en el split 0, idéntico al publicado), lo que valida el código y el protocolo de
evaluación. El modelo trata el video como una secuencia plana de frames y sus números
provienen de corridas individuales, sin varianza entre semillas, que en conjuntos de pocas
decenas de videos cambia la lectura de los resultados.

\section{Extensión A2Summ-X y el diagnóstico contra la fusión}

A2Summ-X, trabajo previo directo de esta tesis, extiende el modelo base hacia una
representación por tomas y jerárquica: adaptadores por modalidad, fusión jerárquica de frames
a tomas y de tomas a video, alineación entre modalidades, compuertas aprendidas por modalidad
y un agregador temporal de clips. Cada componente se activa por separado.

\subsection{Análisis de resultados}
Con 25 corridas por configuración, pruebas t pareadas sobre F1 y métricas de ranking de
Otani \cite{otani2016deep}, se obtuvo que: (i) el texto aporta $\Delta$F1 $= +0.033 \pm
0.062$ ($p = 0.014$) en TVSum, con tendencia positiva no significativa en SumMe; (ii)
reemplazar GoogleNet por DINOv3 o V-JEPA-2 no cambia el resultado ($|\Delta \text{F1}| <
0.001$, $p > 0.8$); y (iii) las mejoras de arquitectura no superan la varianza entre
corridas. El análisis posterior del modelo mostró por qué la señal visual pesaba tan poco: la
alineación inyecta la transcripción en la consulta visual, la compuerta softmax reparte un
peso total fijo entre modalidades (con pesos de texto medidos entre 0.75 y 0.89) y la cabeza
de decisión solo lee la mezcla resultante. Cuatro mediciones independientes coincidieron con
ese diagnóstico: los pesos de la compuerta, el colapso visual de GraphVSum en BLiSS, la
correlación de importancia humana $\rho \approx 0$ fuera de dominio y las secuencias sin
cortes detectables. La conclusión fue que el problema está en el diseño de la fusión y que
cualquier arquitectura que mezcle las representaciones antes de decidir puede reproducirlo.

\section{Segmentación auto-supervisada de escenas: BaSSL}

Mun et al.\ \cite{mun2022bassl} entrenan con tareas contrastivas sensibles a las fronteras
sobre MovieNet \cite{huang2020movienet} y reportan el punto de comparación habitual del área
(F1 47.0 y AP 57.4 en la convención de 318 películas). TranS4mer \cite{islam2023trans4mer} y
MEGA \cite{sadoughi2023mega} refinan la arquitectura temporal y la agregación multimodal;
SCRL \cite{wu2022scrl} y UBoCo \cite{kang2022uboco} exploran la consistencia de escena sin
etiquetas.

\subsection{Análisis de resultados}
Estos sistemas predicen fronteras entre escenas contiguas. No agrupan segmentos del mismo
tema separados en el tiempo ni producen una estructura navegable por tópico. MEGA, el más
cercano a CAVA en el uso de varias modalidades, agrega y destila las modalidades hacia una
representación compartida, que es el diseño cuyo problema se documentó en A2Summ-X. Ninguno
de estos trabajos reporta qué modalidades sostienen cada escena detectada.

\section{Modelos visión--lenguaje para segmentación: Scene-VLM}

Berman et al.\ \cite{berman2025scenevlm} presentan Scene-VLM, un VLM ajustado para segmentar
video en escenas. El modelo recibe frames, la transcripción del diálogo y metadatos
opcionales de tomas consecutivas, y predice las fronteras de forma secuencial con un puntaje
de confianza y una justificación en lenguaje natural. Reportan el mejor resultado publicado
en MovieNet-318 (F1 62.1 y AP 66.8).

\subsection{Análisis de resultados}
El trabajo muestra que un VLM puede emitir juicios semánticos sobre unidades temporales, a
cambio del costo de entrenar y ejecutar un modelo grande sobre cada toma, y su salida sigue
siendo una única secuencia de fronteras. CAVA toma dos decisiones a partir de esto. Primero,
usa el VLM (Qwen2.5-VL \cite{bai2025qwen25vl}) congelado y fuera del camino crítico: sus
descripciones por bloque son una señal más que vota. Segundo, no deja la frontera fina en
manos del VLM, porque el trabajo previo midió que su sensibilidad temporal es baja (invertir o
barajar los frames cambia el score solo entre 0.020 y 0.024).

\section{Capítulos de video y tópicos no contiguos: VidChapters}

Ventura et al.\ \cite{ventura2024vidchapters} construyen VidChapters-7M, un corpus de videos
con capítulos anotados por sus propios autores, y entrenan modelos que generan capítulos a
partir de la transcripción. El problema se formula como predicción de fronteras con títulos.

\subsection{Análisis de resultados}
El corpus muestra que los autores de video declaran estructura temática a gran escala. La
formulación es secuencial (una lista de capítulos contiguos) y solo usa texto: no modela
retornos de tema, no usa imagen ni voz y supone un único nivel de granularidad. Para CAVA,
VidChapters es la referencia externa en el dominio de video web, con los capítulos de autor
como verdad de terreno, y sirve de contraste con un mapa que enlaza segmentos no adyacentes,
usa varias señales y produce dos niveles.

\section{Segmentación temática de texto}

La segmentación temática se estudió primero en texto. TextTiling \cite{hearst1997texttiling}
compara bloques adyacentes por similitud léxica y coloca fronteras en los valles de esa
curva; C99 \cite{choi2000advances} ordena una matriz de similitud entre oraciones y la divide
de forma recursiva.

\subsection{Análisis de resultados}
El score de frontera de CAVA por señal sigue la misma idea que TextTiling: un desplazamiento
entre ventanas vecinas de embeddings, con fronteras en los picos. La diferencia está en que
CAVA aplica ese cálculo por separado a cada señal y no solo a la transcripción, y en que
permite que un segmento vuelva a un tópico anterior. Ambos métodos producen una partición
contigua, lo que deja fuera los retornos de tema que interesan aquí.

\section{Fusión temprana y fusión tardía}

Snoek et al.\ \cite{snoek2005early} comparan, en análisis semántico de video, la fusión
temprana (concatenar los descriptores antes de aprender) con la tardía (aprender por
modalidad y combinar las decisiones). La revisión de Baltrušaitis et al.\
\cite{baltrusaitis2019multimodal} clasifica la fusión en métodos independientes del modelo
(temprana, tardía e híbrida) y métodos basados en el modelo.

\subsection{Análisis de resultados}
En Snoek et al.\ la fusión tardía gana en la mayoría de los conceptos, aunque con
diferencias por concepto, y la revisión de Baltrušaitis et al.\ no establece un ganador
general. CAVA pertenece a la familia de fusión tardía: las señales se combinan solo en el
nivel de las decisiones (cortes y etiquetas), nunca en el de las representaciones. Lo que
agrega frente a la fusión tardía clásica es que la decisión combinada es una partición con
retornos y no una etiqueta por muestra, y que una señal puede abstenerse en los tramos donde
no tiene evidencia.

\section{Ensambles de agrupamientos y matriz de co-asociación}

Strehl y Ghosh \cite{strehl2002cluster} formulan el problema de combinar varias particiones
de los mismos datos en una partición de consenso sin acceder a los descriptores originales.
Fred y Jain \cite{fred2005evidence} proponen la acumulación de evidencia: cada partición vota
por los pares de elementos que pone juntos, los votos se acumulan en una matriz de
co-asociación y la partición final se obtiene agrupando esa matriz con enlace simple o
promedio.

\subsection{Análisis de resultados}
La matriz de votos de CAVA, $V_{ij}$, es una matriz de co-asociación ponderada en la que cada
señal aporta una partición, y su modo de enlace promedio corresponde a la variante de Fred y
Jain que evita el encadenamiento del enlace simple. CAVA adapta este marco de tres maneras:
(i) las particiones de cada señal son temporales y se calculan con su propio percentil; (ii)
una señal sin evidencia en una unidad no vota sobre los pares que la incluyen, de modo que el
consenso se calcula solo sobre los votantes activos; y (iii) cada grupo final conserva la
contribución de cada señal. En la literatura revisada, los ensambles de agrupamientos no se
han aplicado a la segmentación temática de video con señales heterogéneas.

\section{SegLoop: del bucle de destilación al motor de señales}

SegLoop, el sistema previo del grupo, validó un bucle de dos redes para segmentar escenas: un
juez VLM congelado genera pseudoetiquetas con las que se reentrena un segmentador ligero (un
transformer temporal sobre tokens de tomas). En MovieNet-318 obtuvo una base de F1 $45.22 \pm
0.63$ y AP $43.10 \pm 0.62$ y, con preentrenamiento sensible a fronteras y una cabeza
contextual, F1 $51.97 \pm 0.47$ y AP $52.69 \pm 0.43$ en test. La fusión de imagen y texto
se descartó porque no superó a la rama visual.

\subsection{Análisis de resultados}
SegLoop cumple dos papeles en esta tesis. Su resultado negativo de fusión (la familia R2 con
texto no superó a R2) reproduce en segmentación lo observado en A2Summ-X para resumen: al
fusionar, la señal adicional no aporta. Además, su infraestructura está verificada de punta a
punta (detección de cortes con TransNetV2 \cite{soucekl2020transnetv2}, reconocimiento de voz
\cite{radford2023whisper}, descriptores SigLIP \cite{zhai2023siglip} y bloqueo del split de
test). CAVA la usa como motor de extracción de señales a través del contrato
\texttt{frame\_awareness} (puntos de cambio, scores, cortes certificados y tramos con huella)
sin duplicar código.

\section{Modelos visión--lenguaje como jueces}

El uso de modelos grandes como evaluadores automáticos se ha extendido en la evaluación de
texto generado. En video, VideoJudge \cite{videojudge2025} entrena jueces multimodales de 3B y
7B parámetros para puntuar respuestas de comprensión de video.

\subsection{Análisis de resultados}
La experiencia del grupo con VLM como jueces matiza estos resultados: el VLM ordena bien
pares de alternativas (ranking contrafactual casi perfecto) pero da scores absolutos poco
fiables, su sensibilidad temporal es baja y sus veredictos heredan sesgos del entrenamiento.
De ahí salen tres reglas de diseño para CAVA: (i) el VLM genera descripciones por bloque, que
votan como cualquier otra señal, y no veredictos finales; (ii) cuando se usa un juez para
validar (G-enlace), el protocolo es ciego, con controles emparejados y pruebas exactas; y
(iii) ninguna compuerta depende de un único juez no validado.

\section{Aporte de la tesis}

A partir de esta revisión, CAVA reúne cuatro elementos que los trabajos anteriores no
combinan:
\begin{enumerate}
  \item Segmentación temática de video sin fusión de representaciones. Cada señal decide sus
        propias fronteras y etiquetas, y el tópico surge de una matriz de co-asociación con
        quórum \cite{fred2005evidence}, adaptada a particiones temporales y con abstención de
        las señales sin evidencia.
  \item Tópicos no contiguos con procedencia. Los segmentos separados en el tiempo se enlazan
        cuando un quórum de señales los agrupa, y cada grupo declara qué señales lo sostienen
        (acuerdo, soporte y coherencia dejando una unidad fuera), información que una
        representación fusionada no conserva.
  \item Versiones para el consumidor. Del mapa se derivan cinco formatos con presupuestos
        verificables, evaluados en cuatro tipos de contenido.
  \item Un inventario del estado de cada módulo (validado, acotado, parcial, no validado) por
        tipo de contenido, basado en compuertas pre-registradas.
\end{enumerate}
